\chapter{Voxel Geometry and OpenVDB}
\label{ch:voxel_geometry}

Voxel geometry\index{voxel geometry} represents three-dimensional shape on a
regular grid of volume elements, or voxels. Where a mesh stores geometry on
its boundary, a voxel representation samples the volume itself, making
topological queries such as inside/outside tests, boolean operations, and
morphological changes trivial and robust. The cost is memory: a dense uniform
grid grows as the cube of its resolution. Voxel geometry became practical at
high resolution through \emph{sparse} representations, most prominently the
OpenVDB library and its hierarchical tree of fixed-size blocks~\cite{museth2013}.
This chapter develops the dense and sparse representations, signed distance
fields, and the conversions between voxel and mesh form. Mesh
voxelization\index{voxelization} itself is treated in
Section~\ref{sec:mesh_voxelization}.

\section{Voxel Grids}
\label{sec:voxel_grids}

\subsection{1. Overview}
A \textbf{voxel grid} is the 3D analogue of a pixel image: a rectangular
array of cubic cells aligned to an axis-aligned bounding box. Each cell
either records a binary occupancy (occupied or empty) or stores a scalar,
vector, or other attribute sampled at the cell center. Grids support
constant-time random access and simple neighbor iteration, which makes them
the substrate for most volumetric algorithms~\cite{akeninemoller2001,nooruddin2003}.

\subsection{2. Definitions}

\begin{definition}[Voxel]\label{def:voxel_grid}
A \emph{voxel} is a volume element: the cubic cell centered at a grid point.
The grid resolution $(n_x, n_y, n_z)$ and the voxel size $s$ determine the
box that the grid represents.
\end{definition}

\begin{definition}[Binary and Attribute Grids]\label{def:binary_grid}
A \emph{binary grid} stores occupancy per voxel (0 empty, 1 occupied). An
\emph{attribute grid} stores values such as a signed distance, density, or
velocity field.
\end{definition}

\begin{definition}[Conservative Voxelization]\label{def:conservative_vox}
A voxelization is \emph{conservative} if a voxel is occupied whenever any
part of the surface passes through it; otherwise it is \emph{exact} (a voxel
is occupied only when it is pierced by the surface)~\cite{nooruddin2003}.
\end{definition}

\subsection{3. Theory}
A dense grid of resolution $n^3$ costs $n^3$ entries and thus $O(n^3)$ bytes.
The occupied surface of a triangle mesh passes through only $O(n^2)$ voxels,
so for high resolution almost all storage is wasted on empty space. This
observation motivates the sparse schemes of
Section~\ref{sec:sparse_voxel} and the hierarchical trees of
Section~\ref{sec:openvdb}. The resolution must balance aliasing (too coarse)
against memory and computation time (too fine); adaptive sampling that keeps
dense cells near the surface and coarsens empty regions is the resolution
used in production systems~\cite{museth2013}.

\subsection{4. Pseudo code}
\begin{algorithmic}
\Function{IsInside}{$p$, grid, $\mathrm{occ}$}
    \State $c \gets \mathrm{ToCell}(p, \mathrm{grid})$
    \State \Return $\mathrm{occ}[c]$
\EndFunction
\end{algorithmic}

\subsection{5. Parameters Selections}
\begin{itemize}
    \item \textbf{Resolution}: choose $n$ so that the smallest surface feature spans several voxels; alias errors appear below ~3 voxels per feature.
    \item \textbf{Bounds}: an axis-aligned bounding box must contain the whole object; padding avoids boundary artifacts.
    \item \textbf{Data type}: one bit per voxel for occupancy, 8--32 bits per voxel for scalar attributes.
\end{itemize}

\subsection{6. Complexity}
Dense grids: $O(n^3)$ space, $O(1)$ access, $O(n^3)$ build. A surface of $O(n^2)$
cells can be voxelized in $O(n^2)$ time plus $O(n^3)$ for the grid itself.

\subsection{7. Usages}
\begin{itemize}
    \item \textbf{Physics simulation}: grid-based fluids, collision and pressure fields.
    \item \textbf{Medical imaging}: CT and MRI produce data natively as grids.
    \item \textbf{Boolean and morphological operations}: union, intersection, erosion, and dilation are voxel-wise.
\end{itemize}

\subsection{8. Testing methods and Edge cases}
\begin{itemize}
    \item \textbf{Surface vs.\ solid}: verify whether boundary voxels are classified as inside.
    \item \textbf{Resolution sensitivity}: re-run the pipeline at two resolutions and compare the extracted surfaces.
    \item \textbf{Boundary alignment}: voxels exactly on the box boundary must be handled consistently.
\end{itemize}

\subsection{9. References}
\begin{itemize}
    \item Akenine-M\"oller, T. (2001). ``Fast 3D triangle-box overlap testing''. Journal of Graphics Tools.
    \item Nooruddin, F. S., \& Turk, G. (2003). ``Simplification and repair of polygonal models using volumetric techniques''. IEEE TVCG.
    \item Wikipedia: \href{https://en.wikipedia.org/wiki/Voxel}{Voxel}
\end{itemize}

\section{Signed Distance Fields}
\label{sec:sdf_voxel}

\subsection{1. Overview}
A \textbf{signed distance field}\index{signed distance field} (SDF) stores,
at every voxel, the signed distance to the nearest surface, negative inside
and positive outside. The zero set of the field is the surface itself, so
blending, boolean operations, and collision queries reduce to arithmetic on
the sampled field~\cite{osher2003}.

\subsection{2. Definitions}

\begin{definition}[Signed Distance Field]\label{def:sdf_volume}
A \emph{signed distance field} is a function $d:\mathbb{R}^3 \to \mathbb{R}$ with
$|d(x)| = \operatorname{dist}(x, \Gamma)$, where $\Gamma$ is the surface,
and $\operatorname{sgn}(d)$ indicates the side of $\Gamma$.
\end{definition}

\begin{definition}[Narrow Band]\label{def:narrow_band}
A \emph{narrow band} stores $d$ only within $w$ voxels of $\Gamma$ (typically
$w = 3$); outside the band the field is clamped to $\pm w\,s$. This bounds
the stored volume to the neighborhood of the surface~\cite{osher2003,museth2013}.
\end{definition}

\subsection{3. Theory}
An SDF is usually constructed either by computing the distance to a triangle
mesh with a fast sweeping or fast marching method, or by reinitializing an
arbitrary level-set function so that $|\nabla d| = 1$. Once it is a true
signed distance, operations are exact up to the grid resolution: the surface
is extracted as the zero level set, and distance queries are direct voxel
lookups. Level-set evolution updates $d$ with a velocity field, moving the
surface while preserving its implicit character~\cite{osher2003}.

\subsection{4. Pseudo code}
\begin{algorithmic}
\Function{BuildSDF}{$\Gamma$, grid, $w$}
    \State initialize voxels near $\Gamma$ with exact distances
    \State propagate distances outward with fast sweeping
    \State clamp all values to $[{-w\,s}, +w\,s]$
    \State \Return $d$
\EndFunction
\end{algorithmic}

\subsection{5. Parameters Selections}
\begin{itemize}
    \item \textbf{Band width $w$}: $w = 3$ is the standard OpenVDB default; larger bands smooth boolean operations at higher cost.
    \item \textbf{Renormalization}: enforce $|\nabla d| = 1$ periodically so the field stays a true distance.
\end{itemize}

\subsection{6. Complexity}
Fast sweeping computes the distance to a grid of $N$ voxels in $O(N)$ passes
of $O(N)$ work; reinitialization of the narrow band is $O(w\,n^2)$ per sweep
for an $n^3$ grid.

\subsection{7. Usages}
\begin{itemize}
    \item \textbf{Boolean operations}: $\min/\max$ of SDFs implement union/intersection.
    \item \textbf{Collision and proximity}: distance queries become lookup + clamp.
    \item \textbf{Level-set surface evolution} for sculpting and fluid boundaries~\cite{osher2003}.
\end{itemize}

\subsection{8. Testing methods and Edge cases}
\begin{itemize}
    \item \textbf{Accuracy}: compare against analytic distance for a sphere.
    \item \textbf{Sign consistency}: verify the inside/outside convention matches the mesh orientation.
    \item \textbf{Thin features}: features thinner than $2s$ may vanish from the field.
\end{itemize}

\subsection{9. References}
\begin{itemize}
    \item Osher, S., \& Fedkiw, R. (2003). ``Level Set Methods and Dynamic Implicit Surfaces''. Springer.
    \item Wikipedia: \href{https://en.wikipedia.org/wiki/Signed_distance_function}{Signed distance function}
\end{itemize}

\section{Sparse Voxel Representations}
\label{sec:sparse_voxel}

\subsection{1. Overview}
Sparse voxel representations store only occupied cells, typically in a
\textbf{sparse voxel octree}\index{sparse voxel octree} (SVO): a tree that
subdivides the box into octants, recursing only where occupied cells
exist. Hierarchical sparse volumes such as OpenVDB push the idea further,
storing dense fixed-size blocks at the leaves and summarized tiles in the
interior~\cite{museth2013}.

\subsection{2. Definitions}

\begin{definition}[Sparse Voxel Octree]\label{def:svo}
An \emph{SVO} is an octree whose nodes store the geometry and material of the
cells they represent; nodes whose subtree is empty are pruned. Resolution
increases only where the object actually exists.
\end{definition}

\begin{definition}[Tile]\label{def:voxel_tile}
A \emph{tile} is an internal node value that summarizes a uniform region of
the volume (for example a constant SDF value) without expanding it into
children. Tiles are the key to storing large homogeneous regions cheaply~\cite{museth2013}.
\end{definition}

\subsection{3. Theory}
An SVO with $n$ occupied voxels uses $O(n \log n)$ nodes and supports
ray-marching and level-of-detail rendering directly on the tree. OpenVDB's
tree generalizes the octree to arbitrary branching and depth (up to five
levels) and adds \emph{active} states: every node or tile is either active
(stores topology) or inactive (implicitly background), so constant regions
are represented by single tiles instead of millions of voxels~\cite{museth2013}.

\subsection{4. Pseudo code}
\begin{algorithmic}
\Function{Insert}{tree, $c$, value}
    \State descend to the leaf containing cell $c$, expanding empty nodes
    \State set the leaf voxel value and mark the voxel active
    \State \Return tree
\EndFunction
\end{algorithmic}

\subsection{5. Parameters Selections}
\begin{itemize}
    \item \textbf{Branching factor and depth}: octrees use 8; OpenVDB uses 16-child levels with 8$^3$ leaves.
    \item \textbf{Tile threshold}: collapse a subtree into a tile when all children share a value.
\end{itemize}

\subsection{6. Complexity}
Insertion and lookup descend $O(\log n)$ levels (bounded by the fixed tree
depth in OpenVDB). Memory is proportional to the number of active voxels and
tiles, not the bounding-box volume.

\subsection{7. Usages}
\begin{itemize}
    \item \textbf{High-resolution volume rendering}: SVO ray tracing of gigavoxel scenes.
    \item \textbf{Production VFX}: OpenVDB effects in film pipelines~\cite{museth2013}.
\end{itemize}

\subsection{8. Testing methods and Edge cases}
\begin{itemize}
    \item \textbf{Deep trees}: verify behavior at the maximum supported depth.
    \item \textbf{Empty regions}: inserting into a fully empty tree must create only the minimal path.
    \item \textbf{Tile/leaf transitions}: adjacent active and inactive regions must share correct boundaries.
\end{itemize}

\subsection{9. References}
\begin{itemize}
    \item Museth, K. (2013). ``VDB: High-resolution sparse volumes with dynamic topology''. ACM Transactions on Graphics.
    \item Wikipedia: \href{https://en.wikipedia.org/wiki/Sparse_voxel_octree}{Sparse voxel octree}
\end{itemize}

\section{OpenVDB: Hierarchical Sparse Volumes}
\label{sec:openvdb}

\subsection{1. Overview}
\textbf{OpenVDB}\index{OpenVDB} is an open-source C++ library, originated at
DreamWorks and now under the Academy Software Foundation, that implements a
hierarchical, sparse, dynamic topology data structure for volumetric
data~\cite{museth2013}. It is the de facto standard for production voxel
geometry, supporting narrow-band signed distance fields, boolean operations,
smoothing, morphing, and level-set tracking at billion-voxel effective
resolutions.

\subsection{2. Definitions}

\begin{definition}[VDB Tree]\label{def:vdb_tree}
A \emph{VDB tree} has at most five levels: a root node with a hash-mapped set
of children, up to three internal levels, and leaf nodes. Leaves store dense
$8 \times 8 \times 8$ blocks; internal nodes store the states of their
children as fixed-size arrays, and may hold a constant value (a tile) instead
of children~\cite{museth2013}.
\end{definition}

\begin{definition}[Active State]\label{def:vdb_active}
A node, tile, or voxel is \emph{active} if its value is explicitly set and
participates in topology; inactive elements are implicitly equal to the
background value. All tree operations respect and maintain these states.
\end{definition}

\subsection{3. Theory}
Each level of the tree multiplies the coordinate space by 16 in each
dimension, so five levels span coordinates of $2^{19}$ per axis (about
$524{,}288$ voxels) while a leaf stores only 512 voxels. Because interior
levels can be represented by tiles, a large flat region costs the same as a
single voxel. Random access descends at most four pointer hops, and topology
is dynamic: voxels can be inserted or deleted in expected $O(1)$ amortized
time through the hash-mapped root, which makes level-set evolution efficient
even as the interface moves~\cite{museth2013}.

\subsection{4. Pseudo code}
\begin{algorithmic}
\Function{GetValue}{tree, $x$}
    \State $n \gets \Call{root}{tree}$
    \For{level $L = 1 \dots 5$}
        \State $i \gets \operatorname{coord2index}_L(x)$
        \If{$n$ holds a tile} \State \Return tile value
        \ElsIf{$n[i]$ is a child} \State $n \gets n[i]$
        \Else \State \Return background value
        \EndIf
    \EndFor
    \State \Return leaf voxel value at $x$
\EndFunction
\end{algorithmic}

\subsection{5. Parameters Selections}
\begin{itemize}
    \item \textbf{Background value}: choose the "empty" value (e.g., $\pm\infty$ for SDF grids) consistently; it determines inactive semantics.
    \item \textbf{Tree configuration}: leaf log-dimension 3 and internal log-dimension 4 are defaults; larger leaf blocks trade memory for iteration speed.
    \item \textbf{Narrow-band mode}: for SDF grids, band width 3 is standard~\cite{museth2013}.
\end{itemize}

\subsection{6. Complexity}
Access is $O(1)$ (at most five level hops, hash lookup at the root).
Memory is proportional to active voxels and tiles; a 3D SDF of effective
resolution $4096^3$ typically requires only a few megabytes.

\subsection{7. Usages}
\begin{itemize}
    \item \textbf{Visual effects}: level-set sculpting, liquid surfaces, and volumetric fire/smoke in film~\cite{museth2013}.
    \item \textbf{Robotics and simulation}: signed distance collision and meshing pipelines.
    \item \textbf{Medical and geometry processing}: fast boolean operations on scanned geometry.
\end{itemize}

\subsection{8. Testing methods and Edge cases}
\begin{itemize}
    \item \textbf{Bounding}: writing outside the current tree extent must grow the root's hash map correctly.
    \item \textbf{Tile collapse}: coalescing a region must not corrupt neighboring active voxels.
    \item \textbf{Level-set reinitialization}: verify the band stays $|\nabla d| = 1$ after many operations.
\end{itemize}

\subsection{9. References}
\begin{itemize}
    \item Museth, K. (2013). ``VDB: High-resolution sparse volumes with dynamic topology''. ACM Transactions on Graphics.
    \item OpenVDB: \href{https://www.openvdb.org}{https://www.openvdb.org}
    \item Wikipedia: \href{https://en.wikipedia.org/wiki/OpenVDB}{OpenVDB}
\end{itemize}

\section{Voxel-to-Mesh Conversion}
\label{sec:voxel_to_mesh}

\subsection{1. Overview}
Converting a voxel field back into a mesh recovers an explicit surface from
an implicit or occupancy representation. The standard algorithm is
\textbf{marching cubes}\index{marching cubes}, which extracts the zero level
set of an SDF (or the boundary of an occupancy grid) by triangulating each
cell independently~\cite{lorensen1987}.

\subsection{2. Definitions}

\begin{definition}[Marching Cubes]\label{def:marching_cubes_grid}
\emph{Marching cubes} enumerates, for every grid cell, the $2^8$ cases
determined by which of the eight corner values are below the isovalue, and
emits the corresponding polygon from a precomputed case table.
\end{definition}

\begin{definition}[Surface Nets]\label{def:surface_nets}
\emph{Surface nets} place one vertex per crossing cell and connect adjacent
vertices into a quad mesh, relaxing the vertices toward the surface; the
result avoids some of marching cubes' topology artifacts~\cite{gibson1998}.
\end{definition}

\subsection{3. Theory}
Marching cubes produces watertight surfaces for well-sampled SDFs but has
ambiguous cases (15 base cases, 256 with symmetry and inversion) that must be
resolved consistently to avoid cracks. Because the surface is the zero set of
a sampled field, the mesh resolution is bounded by the grid resolution, and
features thinner than one voxel are lost. Surface nets produce smoother
results with fewer triangles at the same resolution~\cite{gibson1998}.

\subsection{4. Pseudo code}
\begin{algorithmic}
\Function{MarchingCubes}{$d$, $\rho$}
    \State $M \gets$ empty mesh
    \For{each cell with corners $c_0 \dots c_7$}
        \State case $\gets \sum_i 2^i\,[d(c_i) < \rho]$
        \State emit triangles of case from the lookup table
    \EndFor
    \State \Return $M$
\EndFunction
\end{algorithmic}

\subsection{5. Parameters Selections}
\begin{itemize}
    \item \textbf{Iso-value $\rho$}: 0 for a signed distance field; 0.5 for a binary grid.
    \item \textbf{Vertex placement}: linear interpolation on edges vs.\ central placement in the cell.
    \item \textbf{Ambiguity resolution}: use the asymptotic-decider criterion for consistent topology.
\end{itemize}

\subsection{6. Complexity}
Marching cubes runs in $O(N)$ time for $N$ cells and outputs $O(N)$
triangles; output is proportional to the surface area, not the volume.

\subsection{7. Usages}
\begin{itemize}
    \item \textbf{Surface extraction} from CT/MRI scans.
    \item \textbf{Mesh generation} from SDF sculpting and level-set simulations.
    \item \textbf{Mesh repair}: voxelize a broken mesh, then extract a clean, watertight one.
\end{itemize}

\subsection{8. Testing methods and Edge cases}
\begin{itemize}
    \item \textbf{Watertightness}: verify no cracks at cell boundaries (ambiguity cases).
    \item \textbf{Normals}: orient triangles consistently from the field gradient.
    \item \textbf{Empty grids}: fields with no sign change must produce no triangles.
\end{itemize}

\subsection{9. References}
\begin{itemize}
    \item Lorensen, W. E., \& Cline, H. E. (1987). ``Marching cubes: A high resolution 3D surface construction algorithm''. Computer Graphics (SIGGRAPH).
    \item Gibson, S. F. F. (1998). ``Constrained elastic surface nets''. MICCAI.
    \item Wikipedia: \href{https://en.wikipedia.org/wiki/Marching_cubes}{Marching cubes}
\end{itemize}

\section{Applications}
\label{sec:voxel_applications}

\begin{itemize}
    \item \textbf{Visual effects}: volumetric fluids, explosions, and organic sculpting in production using OpenVDB~\cite{museth2013}.
    \item \textbf{Autonomous driving and robotics}: occupancy grids and signed-distance collision maps for perception and planning.
    \item \textbf{Medical imaging}: segmentation, surface extraction, and morphometry from volumetric scans~\cite{lorensen1987}.
    \item \textbf{Computational fabrication}: voxel models drive 3D printing and material assignment.
\end{itemize}
